\section{Runtime loop and modules}
\label{sec:runtime}

\subsection{In-game frame profiler}
\symbolref{FrameProfiler} is a main-thread, engine-owned 120-frame timing ring.
Sequential steady-clock marks partition each measured loop body into input,
camera, optional and required updates, scene preparation, asset pumping, command
generation/submission, backend presentation, frame limiting, and remaining time.
The renderer can return the timestamp between submission and presentation;
ordinary untimed calls keep their previous behavior. GPU execution and concurrent
simulation workers are excluded. Disabled sampling performs no clock queries.

DebugModule borrows the collector explicitly and owns a bounded GameVisual pie
with text. F6 or \texttt{profiler on} enables it; keys 1--3 drill into groups and
0 returns to Frame. Values refresh four times per second, including resize and
navigation updates. The console takes priority. Visible profiler navigation
captures both number events and held-key polling before the product update.
FPS and memory visibility remain independent. See
\pathref{docs/frame-profiler.md} and decision D-PROF1.

\subsection{Illumo application runtime}
\pathref{Illumo/Source/Engine/Application.cpp} owns logger lifetime, argument
parsing mechanics, module registration, the frame loop, and the explicit
process exit code. IllumoGame supplies only declarative game metadata and
callbacks through \symbolref{IllumoApplicationDefinition}:
\begin{lstlisting}
int RunIllumoApplication(int argc, char** argv,
                         IllumoApplicationDefinition application) {
  LoggerLifetime logger;
  Illumo illumo(/* application name + environment path */);
  application.applyDefaults(&illumo.environment());
  if (SysCmdLine::ParseCommandLine(argc, argv, &illumo.environment(),
                                   application.commandLine).shouldExit())
    return /* requested exit code */;
  if (!illumo.initialize()) return 1;
  illumo.addModule(application.createRequiredModule(&illumo.environment()), Required);
  illumo.addModule(make_unique<DebugModule>(), Optional); // Debug only
  if (!illumo.startModules()) return 1;
  while (!illumo.shouldClose()) {
    illumo.update(/* steady_clock delta */);
    illumo.render();
  }
  illumo.shutdown();
  return 0;
}
\end{lstlisting}
The engine-owned platform entry obtains the game definition from
\symbolref{CreateIllumoApplication} and invokes this runner. No loop decision
depends directly on GLFW. Help/version and validation may return before
graphics initialization.

The definition in
\pathref{IllumoGame/Source/Game/IllumoGameApplication.cpp} names the product,
describes its CA options/help, supplies its defaults callback, and creates the
initial required module (\symbolref{MainMenuModule} by default, or
\symbolref{CellGameModule} if direct launch was configured). It does not perform system work.

\subsection{Illumo host ownership}
\symbolref{Illumo} owns the long-lived generic services
(\symbolref{EnvVars}, \symbolref{RenderWindow}, \symbolref{Camera},
\symbolref{Renderer}, \symbolref{AssetManager}, \symbolref{CommandRegistry},
\symbolref{CommandLine}, \symbolref{InputManager}, \symbolref{Scene}) and its
registered modules. \symbolref{IllumoConfig} carries only the application name,
configuration-file path, and injectable factories.

Construction loads generic host defaults. The runner then invokes IllumoGame's
CA-default callback and feeds its option descriptors to engine-owned
\symbolref{SysCmdLine} before fallible initialization. The window and backend
factories report failure; Illumo owns exactly one backend initialization attempt
and transfers the resulting \texttt{std::unique\_ptr<IBackend>} to the renderer.
Library code logs and returns failures rather than calling \texttt{std::exit}.

\symbolref{IllumoContext} is a non-owning pointer bag frozen after successful
host initialization and passed to modules at \symbolref{Start}. It exposes
\symbolref{IModuleHost} allowing active modules to request deferred transitions
at frame boundaries. It contains no game validation helpers or product assumptions.

\subsection{IModule contract and failure policy}
\begin{lstlisting}
class IModule {
public:
  virtual bool Start(IllumoContext* context) = 0;
  virtual void Update(double dt) = 0;
  virtual void DispatchDrawables(Scene* scene) = 0;
  virtual void Exit() = 0;
  virtual bool OnCloseRequested() { return true; }
};
\end{lstlisting}

\begin{itemize}
  \item \textbf{Start} binds services and allocates module state.
  \item \textbf{Update} advances simulation and input behavior.
  \item \textbf{DispatchDrawables} contributes this frame's non-owning draw list.
  \item \textbf{Exit} releases module-owned state.
  \item \textbf{OnCloseRequested} may defer native close for product confirmation.
\end{itemize}

The runner invokes \texttt{processCloseRequest} explicitly before termination.
Only started modules participate; acceptance is the default. Deferral clears the
native close flag and continues updating/rendering. IllEd shares its existing
Save/Discard/Cancel flow with native close and records approved discard so a
dirty document does not prompt repeatedly. Required transition failure remains
terminal. Custom windows used with deferring modules implement close-flag
cancellation. The const \texttt{shouldClose} query has no callbacks.

Each registration is \texttt{Required} or \texttt{Optional}. A rejected optional
module is destroyed immediately and never receives update/draw/exit. Required
startup failure rolls back all accepted modules in reverse order and fails
startup. Exceptions from Start or Exit are contained, logged, and do not stop
remaining rollback attempts. At runtime, the active required module can request
transition to another module via \symbolref{IModuleHost::RequestTransition}, which
exits the current module, purges input event queues, clears scene drawables, and
starts the incoming module. The engine runner registers the game-supplied initial
module as required and engine-owned Debug-build \symbolref{DebugModule} as optional.
Optional overlays update before the required product module so console input is
global, then dispatch after it so console, FPS, and renderer-demo drawables sit
on top. Unconsumed key and character events are discarded at the end of each
host update.

\subsection{Frame flow (conceptual)}
\begin{center}
\begin{tikzpicture}[
  font=\scriptsize,
  node distance=0.55cm,
  flowstep/.style={draw, rounded corners, align=center, fill=blue!6, minimum width=4.2cm},
  arr/.style={-{Latex}}
]
  \node[flowstep] (1) {InputManager::update};
  \node[flowstep, below=of 1] (2) {Camera::Update};
  \node[flowstep, below=of 2] (3) {optional overlays: Update(dt)};
  \node[flowstep, below=of 3] (4) {required module: Update(dt)};
  \node[flowstep, below=of 4] (5) {discard leftover key/char events};
  \node[flowstep, below=of 5] (6) {Scene clear + required then overlay draw};
  \node[flowstep, below=of 6] (7) {Renderer::RenderScene (tokens + submit)};
  \node[flowstep, below=of 7] (8) {IBackend::EndFrame (swap)};
  \foreach \a/\b in {1/2,2/3,3/4,4/5,5/6,6/7,7/8} \draw[arr] (\a) -- (\b);
\end{tikzpicture}
\end{center}

\symbolref{RenderWindow} defaults to swap interval one. Persisted
\texttt{vsync=0} selects uncapped profiling, while Debug metrics distinguish
paced swap completions from CPU submissions.

The global F11 shortcut asks the window to toggle fullscreen. The window
publishes the resulting state to EnvVars, and the host reports that value
without changing it again. A rejected toggle preserves the prior setting.

IllMeshViewer consumes queued Press events for O (open), F/R (reset camera),
G (grid), X (wireframe), and B (skybox). Hold/repeat and release events do not
execute these actions. Shortcuts are suppressed while the console is open;
unrelated keys remain available to the host. No shortcut latch is shared
between viewer instances.

IllEd body selection records the initial hit offset on the same constraint plane
used during dragging. Holding or releasing an off-center click without moving
does not translate the object or dirty the document. If the initial drag-plane
intersection fails, selection still succeeds but dragging does not begin.

In 3D, IllEd selects the nearest forward screen-ray intersection with an
object's transformed local bounds. Parent transforms, nonuniform scale and
rotation are included; inherited hidden/disabled state excludes candidates.
Singular or nonfinite bounds are skipped. Bounds follow attachment recipes,
with a small proxy for empty nodes; equal-distance ties prefer later document
entries. This is bounds-based picking, not triangle-exact picking. 2D picking
and ground-plane placement retain their existing paths. A 3D selection ray
need not intersect the ground plane.


\subsection{Process memory diagnostics}
The optional Debug/RelWithDebInfo module composes FPS and memory in one
\texttt{GLString} status panel. F3 and \texttt{fps} control only FPS;
\texttt{memory [on|off|toggle]} controls the persisted \texttt{showMemory}
flag, initially off. Each section may be visible independently. The memory
rows show current working set, process-lifetime peak working set and private
committed memory in MiB with one decimal place. Working set includes shared
pages; private commit need not be resident. These are whole-process values,
not GPU memory or allocation/leak tracking.

The internal presentation state maintains independent refresh clocks. Memory
is queried immediately on enable, then once per second while visible; unchanged
text does not rebuild the drawable. The platform-neutral memory query returns
byte counts and availability. Windows obtains them through the process-memory
API; unsupported scaffolds and failures display ``Memory: unavailable'' and
retry at the normal interval. Native headers stay in the Windows platform
implementation. No worker or new rendering path is introduced.
\subsection{Input registration retirement (2026-09-11)}
InputManager bounds live contexts to 32, reuses released slots and issues distinct
manager-local IDs. Invalid activation preserves selection; active retirement
selects neutral input. CellGameModule releases its registration on Exit and
rejects exhausted startup before allocating domain state. Repeated module entry
must not consume permanent input slots.
